(a) The underlying space is noetherian, and \exref{II.2.9} gives a unique generic point for every nonempty irreducible closed subset.

(b) A minimal nonempty closed subset is irreducible, and every point of it is generic. Uniqueness therefore makes it a singleton.

(c) If two distinct points had the same open neighborhoods, they would have the same closure and would be two generic points of that irreducible closed subset, a contradiction.

(d) An open set missing the generic point has closed complement containing its closure, so it is empty.

(e) A point is minimal for specialization exactly when its closure is a singleton. A point is maximal exactly when its irreducible closure is an irreducible component: inclusion of irreducible closed subsets is equivalent to specialization of their generic points. Every point closure is contained in an irreducible component, so these describe all maximal points. Closed subsets contain the closures of their points and are stable under specialization. Taking complements shows that opens are stable under generization.

(f) Closed subsets of $t(X)$ are precisely the $t(Y)$ for closed $Y\subseteq X$, with inclusions corresponding. Thus $t(X)$ is noetherian. Such a nonempty $t(Y)$ is irreducible exactly when $Y$ is irreducible, and then its unique generic point is the point $Y\in t(X)$, since the closure of $Z\in t(X)$ is $t(Z)$. Hence $t(X)$ is a Zariski space.

If $X$ is a Zariski space, $\alpha:x\mapsto\overline{\{x\}}$ is bijective by existence and uniqueness of generic points, and it identifies each closed subset $Y$ with $t(Y)$. Thus it is a homeomorphism. Conversely, a homeomorphism $X\cong t(X)$ makes $X$ a Zariski space.
